% RESEARCH_DRAFT_SHA256: eff3fced48f934481b4445fbdbdde2ca6585e7d48533fb49c576bbdb9d0d17a7
% Layout and source-reference wrappers expanded for the web renderer.
\documentclass[11pt]{article}
\usepackage[margin=25mm]{geometry}
\usepackage{amsthm}
\usepackage{fontspec,unicode-math}
\setmainfont{Latin Modern Roman}
\setmathfont{Latin Modern Math}
\usepackage{amsmath,amssymb,mathtools,xurl,hyperref}
\hypersetup{colorlinks=true,urlcolor=blue,linkcolor=blue}
\setcounter{secnumdepth}{0}
\setlength{\parindent}{0pt}
\setlength{\parskip}{5pt}
\setlength{\emergencystretch}{3em}
\newcommand{\sref}[2]{\hyperlink{source:#1:#2}{[S#2]}}
\newcommand{\cataloguescope}[1]{\paragraph{Source record.}#1}
\newcommand{\floor}[1]{\left\lfloor#1\right\rfloor}
\newtheorem{proposition}{Proposition}
\newtheorem{lemma}[proposition]{Lemma}
\begin{document}
\paragraph{Research manuscript.} OpenAI \texttt{gpt-6-astra}, reasoning \texttt{ultra}; 11 October 2026.
\paragraph{Preserved source body.} The complete catalogue mathematical text and sources follow with layout and source links expanded, including source qualifications and its historical preparation notes. New research begins at the section headed Research attempt.

\section{Two-Coloring Light Tournaments}\label{statement}

\subsection{Definitions and mathematical statement}

\paragraph{Shared notation.}Standard number systems are $\mathbb N,\mathbb Z,\mathbb Q,\mathbb R,\mathbb C$. Specialized notation and index ranges follow the cited source and the definitions below.


A tournament has exactly one directed arc between every pair of distinct vertices. A vertex set is acyclic if its induced tournament has no directed cycle; two-colorability partitions all vertices into two such sets.
\paragraph{Record scope and qualifications.}
The catalogue title is retained as metadata. Named certificates, contracts, bounds or frameworks in that title are not asserted as results of the original source extract.
The definition of light tournament is not supplied by the question and must be taken from the cited report; it is not replaced by an arbitrary forbidden-subgraph condition.

Source-record limitation: The defining light-tournament condition is absent from the supplied extraction.
\paragraph{Statement recorded in the supplied source.}
Is there a polynomial-time algorithm that partitions every two-colorable light tournament into two acyclic vertex sets, or is deciding two-colorability NP-hard even for light tournaments? \href{https://doi.org/10.4171/owr/2024/50}{[S2]}
\paragraph{Source status.}
The supplied catalogue labels this source record open; that metadata was checked against the supplied snapshot on 2026-10-11. The label is preserved for provenance and is not an independent verification of present open status or mathematical correctness. The catalogue literature assessment is dated 2026-08-22.
\subsection{Short English statement}
Find a polynomial-time two-coloring for the specified light tournaments or prove hardness of deciding two-colorability.
\subsection{Sources}
\begin{itemize}
\item[S1] Ulam AI and the UnsolvedMath Contributors, UnsolvedMath: A Curated Collection of Open Mathematics Problems, catalogue entry. Catalogue curation and metadata: CC BY 4.0. The source extract is retained; definitions, English summaries and qualifications are editorial additions.
\url{https://huggingface.co/datasets/ulamai/UnsolvedMath}.
\item[S2] Combinatorial Optimization (2025), Oberwolfach Reports, DOI 10.4171/owr/2024/50, pp. 2949--3000
\url{https://doi.org/10.4171/owr/2024/50}.
\item[S3] OWR report
\url{https://oa.tib.eu/renate/bitstreams/706cca58-da6b-4d8c-a979-ebf6e92658f0/download}.
\item[S4] Klingelhoefer--Newman
\url{https://epubs.siam.org/doi/abs/10.1137/23M1602127}.
\item Source-specific Oberwolfach Reports licences, including CC BY-SA 4.0 where applicable, are retained; see THIRD\_PARTY\_NOTICES.md.
\end{itemize}
\paragraph{Preparation and review.}
The original source and UnsolvedMath contributors retain attribution for the source extract. Codex (GPT-6; exact serving revision not exposed) assisted with definitions, English summaries, scope notes and formatting. No new human mathematical authorship or independently verified solution is asserted. Catalogue identity checked against the complete supplied registry and current repository. Distinct upstream records, including repeated questions and broad or damaged extracts, are retained for complete coverage of the source CS category. A repeated record is not a new independent conjecture. Complete mathematical-equivalence and current-literature review has not been performed. Preparation and final review follow contribution prompts C and D, with source gaps disclosed. Structural and rendering checks do not establish a new theorem or open status.
\label{end:statement}

\section{Research attempt}
\subsection{Literature, source context and scope}
The full catalogue body, including qualifications and historical source status, is preserved above. The original workshop report was inspected in the locally cached publisher PDF text. This attempt distinguishes the source question from any explicitly restricted model introduced below. Repeated records retain their own source IDs and do not become independent conjectures. Elementary results and attributed literature results are not claimed as new original-target progress.

\subsection{Recovered definition and attempted 2-SAT encoding}
For an arc $u\to v$, the primary report defines $N(u,v)=\{w:v\to w\to u\}$ and calls the arc light when the induced tournament on $N(u,v)$ is acyclic. A tournament is light when every arc is light. We use this supplied primary definition, leaving the qualified catalogue excerpt intact.
A tournament is acyclic exactly when it has no directed triangle. To see this, a shortest directed cycle of length at least four has a chord between its first and third vertices; either orientation produces a shorter directed cycle (a triangle in one direction). Thus a two-coloring is exactly a binary assignment in which every directed triangle is nonmonochromatic.
A tempting route is to encode these constraints as 2-SAT, perhaps with auxiliary Boolean variables. Even the single directed triangle obstructs a direct exact encoding preserving its three color variables.
\begin{proposition}
The ternary nonmonochromatic relation cannot be the projection of the satisfying assignments of any 2-CNF formula.
\end{proposition}
\begin{proof}
Satisfying assignments of a 2-CNF formula are closed under coordinatewise majority of three assignments. If their majority falsified a two-literal clause, at least two assignments would falsify its first literal and at least two would falsify its second literal. Two subsets of size two of a three-element set intersect, so some assignment would falsify both, a contradiction. Unit clauses obey the same property. Projections retain majority closure: choose witnesses for three projected assignments and take the majority including the auxiliary coordinates.
But $001,010,100$ are all nonmonochromatic and their majority $000$ is monochromatic. This contradicts the necessary closure property.
\end{proof}
The directed triangle is light because each arc neighborhood is a singleton. So lightness does not make every individual triangle constraint directly expressible by existential 2-CNF. This does not rule out a more sophisticated polynomial algorithm or a decision reduction that does not preserve all assignments.
\subsection{Verification and remaining gaps}
All eight assignments were checked, giving six legal colorings and the displayed majority violation. The necessary global structure of light tournaments has not been exploited to obtain two colors, and no hardness reduction is proved.

\section{References and source attribution}
The primary Oberwolfach workshop report, \url{https://doi.org/10.4171/owr/2024/50}, was inspected on 11 October 2026. Relevant source passages: 2024/50 p. 2993, definition of light arc and exact two-color question. Catalogue and contributor attribution remain in the original Sources section. Mathematical arguments labeled as this attempt are proved above; cited external results remain attributed to their authors and are not independently reproved or claimed as our discoveries.


\paragraph{Integration review.} Codex (GPT-6; exact serving revision not exposed) reviewed the written arguments and expanded layout and source-reference wrappers. Repository metadata and the production web renderer were checked using Python 3.13 and JavaScript with MathJax 3.2.2. This model-assisted review is not independent human proof verification; no PDF compilation or formal proof checking is asserted.

\section{Completion Estimate}
\noindent\textbf{COMPLETION: 0\%}
No substantiated new progress toward the entire original target is claimed. The worked results and failed route are delimited above; a known result, illustrative model, or auxiliary obstruction is not counted as original-target completion.

\section{Final status and attribution}
\textbf{Attempt status: unresolved.} This describes this research attempt, not a new ruling on the source catalogue's status or on external literature claims.
The mathematics and this self-audit were generated by OpenAI \texttt{gpt-6-astra}, reasoning \texttt{ultra}, on 11 October 2026. Codex (GPT-6; exact serving revision not exposed) orchestrated the batch and contributed formatting/repository review. No human mathematical authorship, independent human proof checking, or formal verification is claimed. The historical preparation paragraphs in the reproduced catalogue body concern that body only.

Contribution prompt A was applied to the complete definition and supplied source record, followed by the adversarial checks of prompt B. Local reproduction files are \texttt{batch011/checks.py}, \texttt{computations.json}, and the per-ID audit JSON. They were not public links at preparation time. Numerical checks are finite and do not prove unrestricted results. No LaTeX engine was found locally; compilation was not performed by this research agent.

\end{document}
